% draft0: 後で使う定義だけを置く。先行研究との比較は次節へ分離する。
\section{背景}
\label{sec:background}

\subsection{メトリクスとサービスレベルRCA}
\begin{itemize}
    \item \textbf{書くこと：} サービス集合，メトリクス集合と所属関係，正常区間・障害後区間，既知の境界時刻，正解サービスと出力順位の記号を定義する．
    \item \textbf{注意：} 異常検知・原因サービス特定・原因メトリクス特定を区別する．正解ラベルと障害種別は評価・事後分析にのみ使う（\hyperref[ev:protocol]{E2}）．
\end{itemize}

\subsection{自己回帰モデルと予測誤差}
\begin{itemize}
    \item \textbf{書くこと：} 切片付きAR($p$)，自己相関，一段先予測と多段再帰予測，定常性の条件，インパルス応答と予測誤差分散を導入する．
    \item \textbf{式案：} $x_t=c+\sum_{j=1}^{p}\phi_jx_{t-j}+\varepsilon_t$．記号と仮定をここで固定し，正常時の学習手順は提案手法へ置く．
    \item \textbf{注意：} companion行列の固有値が単位円内という条件と，AR特性多項式の根が単位円外という条件を混同しない．TODO：定義・予測分散の典拠を確認する．
\end{itemize}

\subsection{Bayes Factorによる分布比較}
\begin{itemize}
    \item \textbf{書くこと：} 周辺尤度，$\mathrm{BF}_{10}=p(D\mid H_1)/p(D\mid H_0)$，平均・分散に対するNormal--Inverse-Gamma（NIG）事前分布，GLRTとの違い．
    \item \textbf{注意：} BFは仮説間の比であり原因サービスの事後確率ではない．推定済みARに基づく誤差系列へ適用する際の仮定は\ref{sec:method-bf}節で明示する．
    \item \textbf{TODO：} 周辺化の導出と典拠を照合する．既知の統計的性質を本研究の新規貢献として扱わない（実装照合：\hyperref[ev:method]{E1}）．
\end{itemize}
